% main.tex
% ============================================================
%  فایل اصلی کتاب «دولت، ملت و تنوع قومی»
% ============================================================

\input{preamble}

\begin{document}

% ---- Title Page ----
\input{titlepage}

% ---- Front Matter ----
\frontmatter
\tableofcontents
\clearpage
\listoffigures
\clearpage
\listoftables

% ---- Preface ----
\chapter*{پیش‌گفتار}
\addcontentsline{toc}{chapter}{پیش‌گفتار}

\begin{quotebox}
«هنر حکمرانی، هنر زیستن با تفاوت‌هاست.»
\hfill — ویل کیملیکا
\end{quotebox}

\vspace{6pt}

این کتاب تلاشی است برای پاسخ به پرسشی بنیادین در سیاست تأسیسی:
چگونه می‌توان ساختاری سیاسی طراحی کرد که تنوع قومی و ملّی را نه تهدید، بلکه فرصت بشمارد؟

در بخش نخست، مبانی نظری شکل‌گیری دولت و رابطه‌ی آن با مفاهیم ملت و قومیت واکاوی می‌شود. بخش دوم، الگوهای سیاسی مدیریت تنوع — از تمرکزگرایی مطلق تا کنفدراسیون — را با نمونه‌های عینی بررسی می‌کند. بخش سوم مطالعات تطبیقی مفصلی از هندوستان، بریتانیا، سوئیس، بلژیک و چند کشور دیگر ارائه می‌دهد. و سرانجام بخش چهارم، به‌طور ویژه به ایران اختصاص یافته و شامل تحلیل وضعیت، چالش‌ها و نقشه‌ی راه اجرایی است.

\simplesep

\vspace{6pt}
امید آنکه این اثر بتواند سهمی هرچند اندک در گفت‌وگوی ملّی درباره‌ی آینده‌ی ساختار سیاسی ایران داشته باشد.

\vspace{1cm}
\begin{flushleft}
\textcolor{PrimaryDark}{\bfseries پژوهشگر}\\
بهار ۱۴۰۴
\end{flushleft}

% ===========================================================
\mainmatter
% ===========================================================

% ---- PART I ----
\partpage{بخش اول}{مبانی نظری و مفهومی}

\input{chapter1}
\input{chapter2}

% ---- PART II ----
\partpage{بخش دوم}{الگوها و مطالعات تطبیقی}

\input{chapter3}
\input{chapter4}

% ---- PART III ----
\partpage{بخش سوم}{ایران: تحلیل و راهبرد}

\input{chapter5}
\input{chapter6}

% ---- Back Matter ----
\backmatter

\input{appendix}

\printbibliography[
  heading=bibintoc,
  title={کتاب‌نامه}
]

\printindex

\end{document}